%=============================================================
%  preamble.tex  ―  パッケージ・設定・マクロ定義
%=============================================================

% ── 日本語フォント（LuaLaTeX） ───────────────────────────────
% luatexja-preset が luatexja-fontspec を内包するので単独でロード
% macOS (Hiragino フォント): hiragino-pron
% Linux/CI (Noto フォント) : noto
\usepackage[hiragino-pron]{luatexja-preset}

% ── 数式 ──────────────────────────────────────────────────────
\usepackage{amsmath}
\usepackage{amssymb}
\usepackage{mathtools}

% ── 図・表 ────────────────────────────────────────────────────
\usepackage{graphicx}
\graphicspath{{figures/}}           % 図の検索パス
\usepackage{booktabs}               % \toprule, \midrule, \bottomrule
\usepackage{threeparttable}         % 表の注釈
\usepackage{longtable}              % 複数ページにまたがる表
\usepackage{multirow}               % セル結合
\usepackage{array}                  % 列幅指定 (>{...}p{...} など)
\usepackage{tabularx}               % 幅指定表
\usepackage{float}                  % [H] 指定
\usepackage{caption}
\captionsetup{
  font       = small,
  labelfont  = bf,
  format     = hang,
  width      = 0.9\linewidth,
}
\captionsetup[table]{position=above}
\captionsetup[figure]{position=below}

% ── 数値・単位 ────────────────────────────────────────────────
\usepackage{siunitx}
\sisetup{
  group-separator   = {,},
  group-minimum-digits = 4,
}

% ── ハイパーリンク ───────────────────────────────────────────
\usepackage[
  colorlinks = true,
  linkcolor  = black,
  citecolor  = black,
  urlcolor   = black,
  bookmarks  = true,
]{hyperref}
\usepackage{url}

% ── 文献 (BibLaTeX) ──────────────────────────────────────────
\usepackage[
  style       = apa,      % APA形式（経済学・経営学の標準）
  backend     = biber,
  natbib      = true,     % \citep, \citet が使える
  sortcites   = true,
  maxcitenames = 2,
  maxbibnames  = 99,
]{biblatex}

% ── カラー・強調 ─────────────────────────────────────────────
\usepackage{xcolor}
\definecolor{myblue}{RGB}{46, 64, 87}
\definecolor{myred}{RGB}{180, 50, 50}

% ── コード表示 ────────────────────────────────────────────────
\usepackage{listings}
\lstset{
  basicstyle   = \ttfamily\small,
  keywordstyle = \color{myblue}\bfseries,
  commentstyle = \color{gray}\itshape,
  stringstyle  = \color{myred},
  numbers      = left,
  numberstyle  = \tiny\color{gray},
  frame        = single,
  breaklines   = true,
  captionpos   = b,
}

% ── その他 ────────────────────────────────────────────────────
\usepackage{setspace}
\setstretch{1.5}           % 行間を 1.5 倍（修論標準: 1.5〜1.6）
\usepackage{enumitem}
\usepackage{pdfpages}      % 外部PDFの取り込み（付録用）

% ── ヘッダ・フッタ ───────────────────────────────────────────
% \usepackage{jlreq-deluxe}  % jlreq用の見出しスタイル拡張（インストール済みなら有効化）
% ヘッダ・フッタのカスタマイズは jlreq の \ModifyHeading 等で行う

%=============================================================
% カスタムマクロ
%=============================================================

% 統計用記号
\newcommand{\AUC}{\mathrm{AUC}}
\newcommand{\SHAP}{\mathrm{SHAP}}
\newcommand{\XGB}{\textrm{XGBoost}}
\newcommand{\Cluster}[1]{\textbf{C#1}}   % \Cluster{0} → C0

% 注釈記号（表内）
\newcommand{\sig}[1]{$^{#1}$}  % \sig{***}

% TODO マーカー（執筆中のみ表示）
\newcommand{\TODO}[1]{\textcolor{red}{\textbf{[TODO: #1]}}}

% 引用ショートカット
% \tcite{key} → 著者 (年) のインライン引用
\newcommand{\tcite}[1]{\citeauthor{#1} (\citeyear{#1})}
